\section{Conclusion}
\label{sec:conclusion}

A fixed-size delta-rule state, the constant-memory endpoint of the
long-context cache trade-off, has a measurable capacity frontier, and on
this testbed that frontier is now located rather than bracketed: recall
collapses in a sharp transition at $K/d_{\mathrm{state}} = 0.5455$ at
$d_{\mathrm{state}} = 64$, % evidence: C2
the same ratio window is loss-free at $d_{\mathrm{state}} = 128$,
% evidence: C3
the $d_{\mathrm{state}} = 80$ frontier excludes ratio invariance,
% evidence: C10
and capacity in raw bindings grows with dimension at an exponent
consistent with state bytes and inconsistent with state width.
% evidence: C15
The most natural single-variable account of the growth, anchor-table
coherence, fails a frozen dose-response test under two injection
structures. % evidence: C6
For provisioning fixed-size memories, the practical rule is that safe
load ratios do not transfer across state sizes. For the science of
fast-weight memory, the residual question is now precisely posed:
something other than the load ratio governs the frontier's location, and
scalar table coherence is not it.
